\subsection[Berechnung des Greifbereichs]{Berechnung des Greifbereichs \myAuthor{Chiara Piccolroaz}} \label{control/range_}
Das Ziel der Berechnung des Greifbereichs des Roboters ist es, das Solver-Team über die maximalen und minimalen x/y-Koordinaten für die Abhol- und Pickup-Positionen der einzelnen Puzzelteile zu informieren.

Nämlich berechnet das Steuerungsteam anhand dieser Koordinaten die Bewegung des Roboters, die dann weiter an das Visualisierungsteam weitergeleitet wird. Sind diese Abhol- und Pickup-Punkte nicht erreichbar, so kann keine inverse Kinematik berechnet werden und keine Informationen an das Visualisierungsteam weitergegeben werden.
\\
\\
Zum ersten wurden mittels der Methoden $checkreachabilityatBase()$ und $getmaxrange()$ in $robot.py$ die minimalen und maximalen Streckungen des Roboterarms berechnet, welche 0 cm und 1433 cm betragen (Figure \ref{fig:control_strecke}).

\begin{figure}[ht]
    \centering
    \includegraphics[width=0.9\textwidth]{strecke.png}
    \caption{Maximale Streckung des Roboterarms}
    \label{fig:control_strecke}
\end{figure}

Mithilfe des Satzes des Pythagoras und trigonometrischen Funktionen wurde dann eine potenzielle maximale Fläche berechnet, innerhalb derer das fertige Puzzle optimal plaziert werden könnte (Figure \ref{fig:control_range}).

\begin{figure}[ht]
    \centering
    \includegraphics[width=0.9\textwidth]{range.png}
    \caption{Potenzielle maximale Fläche des fertigen Puzzles}
    \label{fig:control_range}
\end{figure}

 